\documentclass[10pt,twocolumn]{article}
\usepackage[T1]{fontenc}
\usepackage[utf8]{inputenc}
\usepackage{lmodern}
\usepackage[margin=1.8cm,top=2.4cm,bottom=2cm,columnsep=0.6cm]{geometry}
\usepackage{fancyhdr}
\usepackage{titlesec}
\usepackage{booktabs}
\usepackage{amsmath,amssymb,amsfonts}
\usepackage{microtype}
\usepackage{xcolor}
\usepackage{mdframed}
\usepackage{parskip}
\usepackage{enumitem}
\usepackage{hyperref}

\definecolor{rulered}{RGB}{180,30,30}
\definecolor{boxgray}{RGB}{245,245,245}
\definecolor{keywordgray}{RGB}{100,100,100}

\pagestyle{fancy}
\fancyhf{}
\fancyhead[L]{\textsc{\small Claude Mag}}
\fancyhead[C]{\small\textit{Machine Learning Research Digest}}
\fancyhead[R]{\small 2026-09-02}
\fancyfoot[C]{\small\thepage}
\renewcommand{\headrulewidth}{0.8pt}

\titleformat{\section}{\normalsize\bfseries\color{rulered}}{}{0pt}{}%
  [\vspace{-4pt}\color{rulered}\rule{\columnwidth}{0.4pt}]
\titlespacing{\section}{0pt}{8pt}{4pt}

\hypersetup{colorlinks=true,urlcolor=rulered,linkcolor=black}

\begin{document}
\setlength{\parindent}{0pt}
\setlength{\parskip}{4pt}

{\small\color{keywordgray}\textbf{Keywords:}
  diffusion models \textbullet{}
  reinforcement learning \textbullet{}
  reward alignment \textbullet{}
  path-space gradient}\par\vspace{4pt}

{\LARGE\bfseries\raggedright
  Designing Reinforcement Learning for
  Diffusion Models: A Unified Path-Space View}\par\vspace{4pt}

{\large\itshape\raggedright
  Importance Sampling Between Sampling SDEs Connects All
  Major Diffusion-RL Algorithms Under One Principle}\par\vspace{6pt}

{\small\textbf{By} Yixian Xu, Yuanrui Zhang, Shengjie Luo,
  Liwei Wang, Di He}\par\vspace{2pt}

{\small\color{keywordgray}arXiv:2608.14430 \textbullet{}
  Peking University; ByteDance Seed, August 2026}
\par\vspace{2pt}\noindent{\color{rulered}\rule{\columnwidth}{0.6pt}}\vspace{6pt}

Reinforcement learning fine-tuning aligns diffusion models with
task-specific rewards, but its algorithmic landscape is fragmented:
reverse-trajectory methods, forward-matching methods, and value-gradient
approaches each appear to follow different principles. This paper shows
they all arise from a single path-space importance-sampling derivation
applied to the diffusion SDE. The unification reveals a structured
design space whose axes are value-gradient estimation, weight functions,
and sampling strategy. From this space the authors derive a multi-sample
KDE value-gradient estimator that reduces gradient variance 3.2$\times$
and outperforms prior methods by up to 8.3 pp on human preference tasks.

\begin{mdframed}[backgroundcolor=boxgray,linecolor=rulered,linewidth=1pt,
    innerleftmargin=6pt,innerrightmargin=6pt,
    innertopmargin=6pt,innerbottommargin=6pt]
\textbf{\small AT A GLANCE}\par\vspace{4pt}
\begin{description}[leftmargin=0pt,labelsep=4pt,itemsep=2pt,parsep=0pt,
    font=\small\bfseries]
  \item[Problem:] RL methods for diffusion model alignment are
    algorithmically fragmented with no shared theoretical basis,
    preventing principled comparison or design.
  \item[Why it matters:] A unified framework allows systematic ablation
    and derivation of new, more efficient algorithms.
  \item[Method:] Importance sampling between sampling SDEs on trajectory
    space; derive policy gradient, then its value-gradient form; organize
    algorithm space by estimation, weight, and sampling choices.
  \item[Result:] KDE value-gradient estimator cuts gradient variance
    $3.2\times$ and improves human preference alignment $+8.3$ pp
    on SD3.5-M.
  \item[Evidence:] Ablations over all design-space axes; evaluation on
    image generation (SD3.5-M) and vision-language alignment (Qwen-Image).
\end{description}
\end{mdframed}

\section{The Big Picture}

Diffusion models define a generative process through a forward SDE that
gradually adds noise to data and a reverse SDE that denoises from pure
noise to samples. Fine-tuning for reward alignment means biasing the
reverse process toward high-reward $x_0$ without destroying distributional
coverage. Reinforcement learning provides the right framework: treat each
denoising step as a policy action and optimize expected reward.

But several RL algorithms for diffusion models have been proposed
independently, and their relationships are opaque:
\begin{itemize}[itemsep=1pt,parsep=0pt]
  \item \textbf{DDPO/DPO-Diffusion:} Apply PPO-style clipped probability
    ratios over the full reverse trajectory.
  \item \textbf{AlignProp/ReFL:} Backpropagate reward gradients directly
    through the reverse SDE.
  \item \textbf{Forward-matching methods:} Train on reward-weighted
    forward-noised samples.
\end{itemize}

This paper asks: are these the same algorithm with different
approximations, or genuinely different approaches?

\section{Why Existing Approaches Fall Short}

Without a unifying theory, practitioners must choose among algorithms
without understanding what design choices actually matter. This has
several consequences:

\textbf{Ad hoc hyperparameter selection:} Clipping thresholds, sampling
schedules, and weight functions are set empirically; small changes
produce large performance differences.

\textbf{High variance:} REINFORCE-style estimators applied to the full
reverse trajectory have variance that scales with trajectory length. Prior
methods apply heuristics to reduce variance without theoretical grounding.

\textbf{Inaccessible algorithm space:} It is unclear whether a new
estimator can be derived within this framework or whether improvements
require fundamentally new ideas.

\section{Core Mechanism}

The paper applies importance sampling between the current reverse-SDE
measure $p_\theta(\tau)$ and the reward-weighted measure
$p^*(\tau) \propto p_\theta(\tau) R(x_0)$ on the full trajectory space
$\tau = (x_T, \ldots, x_0)$.

The importance weight is:
\[
  w(\tau) = \frac{p^*(\tau)}{p_\theta(\tau)} = \frac{R(x_0)}{\mathcal{Z}}
\]
where $\mathcal{Z} = \mathbb{E}[R(x_0)]$ normalizes the measure.

The policy gradient $\nabla_\theta J(\theta) = \nabla_\theta
\mathbb{E}_{p_\theta}[R(x_0)]$ becomes:
\[
  \nabla_\theta J = \mathbb{E}_{p_\theta}\!\left[R(x_0)
    \nabla_\theta \log p_\theta(\tau)\right]
\]
This is a standard policy gradient on trajectory space. Factoring
$\log p_\theta(\tau) = \sum_t \log p_\theta(x_{t-1} | x_t)$ shows that
each reverse step contributes a per-step gradient weighted by the
downstream reward.

\section{Technical Formulation}

The \textbf{value function} at denoising step $t$ is:
\[
  V^\pi(x_t) = \mathbb{E}_{p_\theta(\tau_{t:0} | x_t)}[R(x_0)]
\]
The policy gradient can be re-expressed in value-gradient form:
\[
  \nabla_\theta J = \mathbb{E}_{x_T}\!\left[\nabla_\theta V^\pi(x_T)\right]
\]
where $x_T \sim \mathcal{N}(0, I)$ is the initial noise.

The paper proposes a \textbf{multi-sample KDE value estimator}:
\[
  \hat{V}(x_t) = \frac{\sum_{i=1}^K R(x_0^{(i)}) K_h(x_0 - x_0^{(i)})}
                       {\sum_{i=1}^K K_h(x_0 - x_0^{(i)})}
\]
where $K$ rollouts $\{x_0^{(i)}\}$ are drawn from $p_\theta(\cdot | x_t)$
and a Gaussian kernel $K_h$ smooths the estimate. Variance is:
\[
  \text{Var}(\hat{V}_\text{KDE}) = O(1/K)
\]
compared to $O(1)$ for the single-sample REINFORCE baseline.

The three design axes of the unified space are:
\begin{enumerate}[itemsep=2pt,parsep=0pt]
  \item \textbf{Weight function} $f(w(\tau))$: identity (vanilla PG),
    clip ($\epsilon$-PPO), binary advantage ($\mathbb{1}[R > c]$)
  \item \textbf{Value-gradient estimation}: direct backprop through SDE,
    REINFORCE, multi-sample KDE, separate critic
  \item \textbf{Sampling strategy}: on-policy, off-policy with
    importance correction, mixed
\end{enumerate}
Existing algorithms correspond to specific combinations; new ones emerge
from unexplored combinations.

\section{Learning or Inference Procedure}

\textbf{Training loop:}
\begin{enumerate}[itemsep=2pt,parsep=0pt]
  \item Sample $M$ initial noises $\{x_T^{(m)}\}$
  \item For each $x_T^{(m)}$, draw $K$ rollouts of the full
    reverse SDE to obtain $\{x_0^{(m,k)}\}$
  \item Compute reward $R(x_0^{(m,k)})$ for each sample
  \item Compute KDE value estimates $\hat{V}(x_t^{(m)})$ at each step
  \item Compute the value gradient and apply an optimizer step to $\theta$
\end{enumerate}

\textbf{Inference:} Standard diffusion model sampling; no change to the
inference pipeline.

\section{What the Guarantee Says}

Under smoothness of $R$ and the reverse SDE transition kernels, the
KDE estimator satisfies:
\[
  \text{Var}(\hat{V}_\text{KDE}) \leq \frac{C \|R\|_\infty^2}{K h^d}
\]
where $d$ is the data dimension and $h$ the kernel bandwidth. Bias
is $O(h^2)$. The optimal bandwidth balancing bias and variance gives
$h = O(K^{-1/(d+4)})$, and the total mean-squared error decays as
$O(K^{-4/(d+4)})$.

\section{Experimental Findings}

\begin{itemize}[itemsep=2pt,parsep=0pt]
  \item \textbf{SD3.5-M (image generation):} KDE estimator outperforms
    single-sample REINFORCE by $+8.3$ pp on human preference alignment;
    outperforms AlignProp by $+4.6$ pp.
  \item \textbf{Qwen-Image (VQA alignment):} $+4.1$ pp absolute over
    AlignProp on target reward; $+2.3$ pp over DDPO.
  \item \textbf{Variance reduction:} KDE with $K=8$ reduces gradient
    variance $3.2\times$ vs.\ single-sample; variance continues
    decreasing to $K=32$ with diminishing returns.
  \item \textbf{Design space ablations:} Binary advantage weighting
    degrades performance $-3.1$ pp vs.\ identity; off-policy with
    importance correction recovers $+1.8$ pp of the gap.
\end{itemize}

\section{Ablations and Interpretation}

\textbf{Number of rollouts $K$:} The best pass-at-1 reward improves
monotonically with $K$ up to $K=16$; wall-clock time scales linearly.
At $K=8$ the compute-performance trade-off is near-optimal.

\textbf{Weight function:} Identity weighting (vanilla PG) outperforms
PPO clipping in the low-data regime; clipping helps stability when
off-policy data is reused.

\textbf{Trajectory length:} Methods that apply the gradient estimator
only at the final denoising step (equivalent to treating the full
reverse trajectory as a black box) underperform full-trajectory
methods by $-5.2$ pp, confirming that per-step value estimates matter.

\textbf{Softmax saturation (for discrete diffusion variants):} The KDE
estimator generalizes to discrete diffusion language models with a
categorical kernel; variance reduction persists ($2.8\times$).

\section*{Reference}

Yixian Xu, Yuanrui Zhang, Shengjie Luo, Liwei Wang, Di He.
\textbf{Designing Reinforcement Learning for Diffusion Models:
A Unified Path-Space View.}
arXiv:2608.14430, Peking University; ByteDance Seed, August 2026.
\url{https://arxiv.org/abs/2608.14430}

\end{document}
